\section{Matrix-valued functions of one real variable}

Matrix-valued functions of a real parameter are differentiated entry by
entry, without fixing a matrix norm.

\begin{lemma}[A first-order quotient rule]
    \label{lem:glm_l_path_quotient_rule}
    \lean{Matrix.hasDerivAt_mul_entries,Matrix.hasDerivAt_trace_of_entries,
        Matrix.differentiableAt_det_of_entries,
        Matrix.differentiableAt_inv_of_entries,
        Matrix.hasDerivAt_entries_of_eventually_eq_mul}
    \leanok
    Let $R(\gamma)$, $L(\gamma)$ and $P(\gamma)$ be finite complex matrices
    depending on a real parameter, with $P$ square. Suppose $R=LP$ on a
    neighbourhood of $\gamma_0$ and $P(\gamma_0)=1$. Suppose also that the
    entries of $R$ and $P$ are differentiable at $\gamma_0$ with derivatives
    $R'$ and $P'$. Then the entries of $L$ are differentiable at $\gamma_0$,
    with derivative $R'-L(\gamma_0)P'$.
\end{lemma}
\begin{proof}
    \leanok
    The determinant of $P$ is a polynomial in its entries. It is therefore
    differentiable at $\gamma_0$, and it is nonzero near $\gamma_0$ because
    $\det P(\gamma_0)=1$. By Cramer's rule the entries of $P^{-1}$ are
    differentiable at $\gamma_0$. Near $\gamma_0$ we have $L=RP^{-1}$, so
    $L$ is differentiable there. The product rule applied to $R=LP$ gives
    $R'=L'P(\gamma_0)+L(\gamma_0)P'=L'+L(\gamma_0)P'$.
\end{proof}

\begin{lemma}[A frame whose derivative stays in its span]
    \label{lem:glm_l_path_frame}
    \lean{Matrix.hasDerivAt_inv_entries,
        Matrix.eventually_exists_eq_sum_smul_of_hasDerivAt,
        Matrix.exists_eq_sum_smul_of_hasDerivAt}
    \leanok
    Let $U\subseteq\mathbb R$ be open and connected, and let $V^i(\gamma)$,
    $i=1,\ldots,M$, be complex matrices of a fixed size depending on
    $\gamma\in U$. Suppose that at every $\gamma\in U$ the entries of each
    $V^i$ are differentiable, with
    \[
        \frac{d}{d\gamma}V^i(\gamma)=\sum_{i'}X(\gamma)_{ii'}V^{i'}(\gamma)
    \]
    for some square matrix $X(\gamma)$, and that there are matrices
    $\widehat V^k(\gamma)$ and a scalar $c(\gamma)$ with
    $c(\gamma)\operatorname{tr}(V^i(\gamma)\widehat V^k(\gamma))=\delta_{ik}$.
    Then for $\gamma_0,\gamma\in U$ there is an invertible matrix $G$ with
    $V^i(\gamma)=\sum_kG_{ik}V^k(\gamma_0)$ for every $i$. No continuity of
    $\gamma\mapsto X(\gamma)$ is assumed.
\end{lemma}
\begin{proof}
    \leanok
    \uses{lem:glm_l_path_quotient_rule}
    Fix $t\in U$ and put
    $Q(\gamma)_{ik}=c(t)\operatorname{tr}(V^i(\gamma)\widehat V^k(t))$. Then
    $Q(t)=1$ and $Q'=XQ$ wherever the derivative hypothesis holds. The
    determinant of $Q$ is differentiable at $t$, so $Q$ is invertible on a
    ball around $t$. On that ball the entries of $Q^{-1}$ are differentiable
    by Cramer's rule, and differentiating $Q^{-1}Q=1$ gives
    $(Q^{-1})'=-Q^{-1}Q'Q^{-1}=-Q^{-1}X$. Hence
    $(Q^{-1}V)'=-Q^{-1}XV+Q^{-1}XV=0$, so $Q^{-1}V$ is constant on the ball
    and equals $V(t)$. Thus $V(\gamma)=Q(\gamma)V(t)$ near $t$.

    The relation ``$V(\gamma)=GV(\gamma')$ for an invertible $G$'' is
    reflexive, symmetric and transitive, and by the first step it holds for
    all $\gamma$ near any $\gamma'\in U$. Its classes are therefore open, and
    since $U$ is connected there is only one class.
\end{proof}

\begin{lemma}[Transport by a continuous matrix differential equation]
    \label{lem:glm_l_path_matrix_ode}
    \lean{Matrix.eqOn_of_hasDerivAt_mul,
        Matrix.eq_mul_initial_of_hasDerivAt_mul,
        Matrix.isUnit_of_hasDerivAt_mul,
        Matrix.eq_mul_initial_of_hasDerivAt_mul_entries,
        Matrix.eq_mul_initial_mul_inv_of_hasDerivAt,
        Matrix.eq_kronecker_mul_initial_mul_inv_of_hasDerivAt}
    \leanok
    Let $I=[\alpha,\beta]$ and $t_0\in I$. A continuous square matrix
    coefficient $X$ has at most one differentiable solution of $F'=XF$
    with prescribed value at $t_0$. If $g'=Xg$ and $g(t_0)=1$, then $g(t)$
    is invertible and $F(t)=g(t)F(t_0)$, also for rectangular $F$.
    For the two-sided rectangular case, let $X:I\to\mathbb C^{n\times n}$
    and $Y:I\to\mathbb C^{m\times m}$ be continuous, and let
    $F:I\to\mathbb C^{n\times m}$, $G:I\to\mathbb C^{n\times n}$ and
    $H:I\to\mathbb C^{m\times m}$ satisfy
    \[
        F'=XF-FY,\qquad G'=XG,\qquad H'=YH,\qquad G(t_0)=1,\quad H(t_0)=1.
    \]
    Here the functions are defined in real neighborhoods of the interval so
    that the derivatives, entrywise for $F$, exist at every point of $I$,
    including its endpoints. Then $G(t)$ and $H(t)$ are invertible and
    \[
        F(t)=G(t)F(t_0)H(t)^{-1}\qquad(t\in I).
    \]
    More generally, let $X_1,X_2,X_3$ be continuous matrix coefficients
    and let $g_\nu'=X_\nu g_\nu$, $g_\nu(t_0)=1$. If
    $L'= (X_1\otimes1+1\otimes X_2)L-LX_3$, then
    \begin{align}
        L(t) & =(g_1(t)\otimes g_2(t))L(t_0)g_3(t)^{-1}.
        \label{eq:glm_l_path_integrated}
    \end{align}
    All derivatives are assumed to exist at every point of $I$.
\end{lemma}
\begin{proof}
    \leanok
    Continuity bounds $\|X(t)\|$ on the compact interval. Uniqueness for a
    uniformly Lipschitz differential equation applies both forwards and
    backwards from $t_0$. Comparing $F$ with $gF(t_0)$ gives transport.
    If $g(t)U=g(t)V$, uniqueness backwards to $t_0$ gives $U=V$;
    finite dimensionality therefore makes $g(t)$ invertible.
    Apply transport columnwise to rectangular matrices. For the two-sided
    equation, the product rule gives
    \[
        (FH)'=(XF-FY)H+F(YH)=X(FH),\qquad (FH)(t_0)=F(t_0).
    \]
    Rectangular transport yields $F(t)H(t)=G(t)F(t_0)$; multiplying by
    the inverse of $H(t)$ proves the two-sided formula. The product rule
    gives $(Lg_3)'=(X_1\otimes1+1\otimes X_2)Lg_3$ and
    $(g_1\otimes g_2)'=(X_1\otimes1+1\otimes X_2)(g_1\otimes g_2)$.
    Transport of $Lg_3$, followed by multiplication by $g_3^{-1}$,
    proves~(\ref{eq:glm_l_path_integrated}).
\end{proof}
